% Part: first-order-logic
% Chapter: proof-systems
% Section: axiomatic-deduction

\documentclass[../../../include/open-logic-section]{subfiles}

\begin{document}

\iftag{FOL}
      {\olfileid{fol}{prf}{axd}}
      {\olfileid{pl}{prf}{axd}}

\olsection{Axiomatische \usetoken{P}{derivation}}

Axiomatische !!{derivation}s zijn de oudste en eenvoudigste logische
!!{derivation}ssystemen. Hun !!{derivation}s zijn eenvoudigweg rijen
!!{sentence}s. Een rij !!{sentence}s geldt als een correcte
!!{derivation} als elke !!{sentence}~$!A$ erin aan een van de volgende
voorwaarden voldoet:
\begin{enumerate}
\item $!A$ is een axioma, of
\item $!A$ is !!a{element} van een gegeven verzameling~$\Gamma$ van !!{sentence}s, of
\item $!A$ is gerechtvaardigd door een afleidingsregel.
\end{enumerate}
Om een axioma te zijn, moet $!A$ de vorm hebben van een van een aantal
vaste !!{sentence}sschema's. Er bestaan veel verzamelingen
axiomaschema's die voor de eerste-ordelogica een deugdelijk, dus correct
en volledig, !!{derivation}ssysteem opleveren. Sommige zijn geordend
naar de connectieven waarop ze betrekking hebben. De schema's
\[
!A \lif (!B \lif !A) \qquad !B \lif (!B \lor !C) \qquad (!B \land !C) \lif !B
\]
zijn bijvoorbeeld gangbare axioma's voor respectievelijk $\lif$,
$\lor$ en~$\land$. Sommige axiomasystemen streven naar een minimaal
aantal axioma's. Afhankelijk van de connectieven die als primitief
worden aangenomen, bestaan er zelfs axiomasystemen met slechts één
axioma.

Een afleidingsregel is een voorwaardelijke uitspraak die een voldoende
voorwaarde geeft waaronder !!a{sentence} in !!a{derivation}
gerechtvaardigd is. Modus ponens is een zeer gangbare regel van dit
soort: als $!A$ en $!A \lif !B$ al gerechtvaardigd zijn, dan is ook
$!B$ gerechtvaardigd. Een stap in !!a{derivation} die de
!!{sentence}~$!B$ bevat, is dus gerechtvaardigd mits zowel $!A$ als
$!A \lif !B$ (voor een zekere !!{sentence}~$!A$) eerder dan~$!B$ in de
!!{derivation} voorkomen.

De op axiomatische !!{derivation}s gebaseerde relatie $\Proves$ wordt
als volgt gedefinieerd: $\Gamma \Proves !A$ dan en slechts dan als er
!!a{derivation} bestaat met de !!{sentence}~$!A$ als laatste formule
(waarbij $\Gamma$ wordt genomen als de verzameling !!{sentence}s in die
!!{derivation} die volgens~(2) hierboven zijn gerechtvaardigd). $!A$
is een stelling als er !!a{derivation} van~$!A$ bestaat waarvoor
$\Gamma$ leeg is, dat wil zeggen als elke !!{sentence} in de
!!{derivation} volgens (1) of~(3) is gerechtvaardigd. De volgende
!!a{derivation} laat bijvoorbeeld zien dat $\Proves !A  \lif (!B \lif
(!B \lor !A))$:
\begin{derivation}
  1. & $!B \lif (!B \lor !A)$ \\
  2. & $(!B \lif (!B \lor !A)) \lif (!A  \lif (!B \lif (!B \lor !A)))$\\
  3. & $!A  \lif (!B \lif (!B \lor !A))$
\end{derivation}
De !!{sentence} op regel~1 heeft de vorm van het axioma $!A \lif (!A
\lor !B)$ (waarbij de rollen van $!A$ en $!B$ zijn verwisseld). De zin
op regel~2 heeft de vorm van het axioma $!A \lif (!B \lif !A)$. Beide
regels zijn dus gerechtvaardigd. Regel~3 is door modus ponens
gerechtvaardigd: als we de formule op die regel afkorten als $!D$, dan
heeft regel~2 de vorm $!C \lif !D$, waarbij $!C$ gelijk is aan $!B \lif
(!B \lor !A)$, dus aan regel~1.

Een verzameling $\Gamma$ is inconsistent als $\Gamma \Proves \lfalse$.
In een volledig axiomasysteem is bovendien $\lfalse \lif !A$ voor
elke~$!A$ bewijsbaar. Als $\Gamma$ inconsistent is, dan geldt dus
$\Gamma \Proves !A$ voor elke~$!A$.

Gottlob Frege gaf in zijn \emph{Begriffsschrift} uit 1879 de eerste
systemen van axiomatische !!{derivation}s voor de logica. Daarom wordt
dit werk vaak beschouwd als het eerste werk van de moderne logica.
Alfred North Whitehead en Bertrand Russell vervolmaakten zulke systemen
in hun \emph{Principia Mathematica}; ook David Hilbert en zijn leerlingen
deden dat in de jaren twintig van de twintigste eeuw. Daarom worden ze
vaak ``Frege-systemen'' of ``Hilbert-systemen'' genoemd. Ze zijn zeer
veelzijdig, doordat voor een logica vaak gemakkelijk een axiomatisch
systeem kan worden gevonden. Omdat de !!{derivation}s zeer eenvoudig
zijn opgebouwd en de systemen slechts een of twee afleidingsregels
hebben, zijn ook uitspraken \emph{over} die !!{derivation}s
betrekkelijk gemakkelijk te bewijzen. In de praktijk zijn ze echter
zeer moeilijk te gebruiken: bewijzen zijn moeilijk te vinden en uit
te schrijven.
\end{document}
